\documentclass[../../../include/open-logic-section]{subfiles}

\begin{document}

\olfileid{sth}{choice}{countablechoice}
\olsection{د شمېر وړ انتخاب}

بې له انتخاب يا ښه ترتيب څخه دا په اسانۍ ثابتېږي چې:

\begin{lem}[په $\Zminus$ کښې]
هر متناهي سټ د انتخاب تابعه لري۔
\end{lem}

\begin{proof}
$a = \{b_1, \ldots, b_n\}$ ونیسئ۔ د اسانتيا
دپاره فرض کړئ چې هر $b_i
\neq \emptyset$ وي۔ نو داسې څيزونه
$c_1, \ldots, c_n$ شته چې
$c_1 \in b_1, \ldots, c_n \in b_n$ وي۔
اوس د
\olref[z][pairs]{prop:pairsconsequences} له مخې
سټ $\{\langle b_1,
c_1\rangle , \ldots, \langle b_n, c_n\rangle\}$
شتون لري؛ او دا د~$a$ دپاره د انتخاب تابعه ده۔
\end{proof}

خو د نامتناهي سټونو په پام کښې نيولو سره چارې
سمدستي لا مبهمېږي۔ د بېلګې په توګه، د پورته
پايلې دا «تر ټولو کوچنے» پراخول وګورئ:

\begin{defish}
\emph{د شمېر وړ انتخاب۔} هر \emph{د شمېر وړ}
سټ د انتخاب تابعه لري۔
\end{defish}

دا د انتخاب د اصل يو ځانګړے حالت دے۔ او څرګندېږي
چې دا اصل ډېر ځله بې له دې کارول شوے چې کاروونکي
يې په ښکاره پوه شوي وي۔ دوه ښې بېلګې يې دلته دي۔
\footnote{د \citet[\S9.4]{Potter2004} او لوکا
انکورواتي له مخې۔}

\begin{ex}
دا يو طبيعي فکر دے: د هر سټ $A$ دپاره يا
$\cardle{\omega}{A}$ وي، يا
$\cardeq{A}{n}$ د کوم $n \in \omega$
دپاره وي۔ دا د دې وجداني مفکورې د بيان يوه
لاره ده چې هر سټ يا متناهي دے يا نامتناهي۔
کانتور او ډېرو نورو رياضيپوهانو دا دعوه بې له
ثبوته کړې وه۔ موږ په احتياط سره دا په
\olref[cardinals][classing]{generalinfinitycharacter}
کښې ثابته کړه۔ خو په هغه ثبوت کښې مو په
$\ZFC$ کښې کار کاوه، ځکه دا مو فرض کړې
وه چې هر سټ $A$ ښه ترتيبېږي، او ځکه
$\card{A}$ خامخا شتون لري۔ يعنې موږ
انتخاب په ښکاره فرض کړے و۔

په حقيقت کښې \citet{Dedekind1888} د
دې دعوې خپل ثبوت په لاندې ډول وړاندې کړ:

\begin{thm}[په $\Zminus + \text{د شمېر وړ انتخاب}$ کښې]
د هر $A$ دپاره يا $\cardle{\omega}{A}$
وي يا $\cardeq{A}{n}$ د کوم
$n \in \omega$ دپاره وي۔
\end{thm}

\begin{proof}
فرض کړئ چې $\cardneq{A}{n}$ د هر
$n \in \omega$ دپاره وي۔ نو په ځانګړي
ډول د هر $n < \omega$ دپاره داسې
فرعي سټ $A_n \subseteq A$ شته چې
عين $\cardexpo{2}{n}$ غړي لري۔ د دې
لړۍ $A_0, A_1, A_2,
\ldots$ په کارولو د هر $n$ دپاره دا
تعريفوو:
\[
	B_n = A_n \setminus \bigcup_{i < n} A_i.
\]
اوس لاندې ته پام وکړئ:
\begin{align*}
	\card{\bigcup_{i < n}A_n} 
	&\leq \card{A_0} + \card{A_1} + \ldots + \card{A_{n-1}}\\
	&=1 + 2 + \ldots + 2^{n-1}\\
	& = 2^n - 1\\
	& < 2^n = \card{A_n}
\end{align*}
\textbf{د سرچينې يادښت (OLSTH-027):} د چاپي
نامساوات په لومړۍ کرښه کښې د اتحاد شاخص
د مخکينۍ تعريف شوې لړۍ له شاخص سره نه
لګېږي؛ د همدې چاپي بڼې له مخې نامساوات
نه ثابتېږي۔ فورمول هماغسې ساتل شوے او
تشه لا نه ده رغول شوې۔
نو هر $B_n$ لږ تر لږه يو غړے،
$c_n$، لري۔ سربېره پردې، $B_n$-ګانې
دوه په دوه بېلې دي؛ نو که $c_n = c_m$
وي، بيا $n = m$ وي۔ خو هر
$c_n
\in A$ دے۔ نو تابعه $f(n) = c_n$
له $\omega \to A$ څخه يوه يو پر يو
تابعه ده۔
\end{proof}
\noindent
دېدېکيند په ښکاره ونه ويل چې د شمېر وړ
انتخاب يې کارولے دے۔ خو \emph{تاسو} يې
کارونه وليدله؟ بيا يې وګورئ۔ (په رښتيا:
\emph{بيا يې وګورئ}۔)

په ثبوت کښې د شمېر وړ انتخاب دوه ځله
کارول شوے۔ يو ځل مو د سټونو لړۍ
$A_0$, $A_1$, $A_2$, \dots\@
ترلاسه کولو ته وکاراوه۔ بيا مو د غړو
$c_n$ د هر~$B_n$ څخه ټاکلو ته وکاراوه۔
سربېره پردې، دلته د انتخاب کارونه نه شي
لرې کېدے۔ \citet[p.~138]{Cohen1966}
ثابته کړې چې دا پايله بې له هر ډول
انتخابه ناکامېږي۔ يعنې له $\ZF$ سره
دا سازګاره ده چې ځينې سټونه له~$\omega$
سره \emph{نه پرتله کېږي}۔
\end{ex}

\begin{ex}
په \citeyear{Cantor1878} کښې کانتور وويل
چې د شمېر وړ سټونو د شمېر وړ اتحاد هم
د شمېر وړ دے۔ هغه ثبوت نه و وړاندې کړے؛
ښايي ثبوت يې ښکاره ګڼلے و۔ موږ بيا په
احتياط سره د دې پايلې يوه عامه بڼه په
\olref[card-arithmetic][simp]{kappaunionkappasize}
کښې ثابته کړې۔ خو زموږ ثبوت انتخاب
په ښکاره فرض کړے و۔ ان د لږې عامې
پايلې ثبوت هم د شمېر وړ انتخاب غواړي۔

\begin{thm}[په $\Zminus + \text{د شمېر وړ انتخاب}$ کښې]
که $A_n$ د هر $n \in \omega$
دپاره د شمېر وړ وي، نو $\bigcup_{n <
\omega} A_n$ هم د شمېر وړ دے۔
\end{thm}

\begin{proof}
بې له عموميت څخه د څه له لاسه ورکولو،
فرض کړئ چې هر $A_n \neq \emptyset$
وي۔ نو د هر $n \in \omega$ دپاره
!!a{surjection} $f_n \colon \omega
\to A_n$ شته۔ تابعه $f \colon \omega
\times \omega \to \bigcup_{n <
\omega} A_n$ د $f(m, n) = f_n(m)$
په وسيله تعريف کړئ۔ پايله له دې راځي چې
$\omega
\times \omega$ د شمېر وړ دے
(\olref[sfr][siz][zigzag]{natsquaredenumerable})
او $f$ يوه !!a{surjection} ده۔
\end{proof}
\noindent
ايا تاسو د شمېر وړ انتخاب کارونه وليدله؟
دا د تابعو لړۍ $f_0$, $f_1$, $f_2$,
\dots\footnote{د انتخاب ورته کارونه په
\olref[card-arithmetic][simp]{kappaunionkappasize}
کښې هم شوې وه، چې هلته مو وويل «د هر
$\beta \in \cardfont{a}$ دپاره يوه
!!a{injection} $f_\beta$ وټاکئ»۔}
د ټاکلو دپاره کارېږي۔ بيا هم، دا پايله
بې له هر ډول انتخابي اصل څخه ناکامېږي۔
په ځانګړي ډول،
\citet{FefermanLevy1963} ثابته کړې
چې له $\ZF$ سره دا سازګاره ده چې
د شمېر وړ سټونو د شمېر وړ اتحاد اصلي
شمېر~$\beth_1$ ولري۔ خو د همدې خبرې
يوه ډېره خوندوره بېلګه د رسل له قوله
دا ده:
\begin{quote}
  د دې بېلګه هغه ميليونر دے چې هر ځل
  به يې د بوټانو له جوړې سره د جرابو
  يوه جوړه هم اخيسته، او له دې پرته يې
  هېڅکله جرابې نه اخيستې۔ د دواړو له
  اخيستلو سره يې دومره مينه وه چې
  بالاخره يې د بوټانو $\aleph_0$ جوړې
  او د جرابو $\aleph_0$ جوړې لرلې\dots\@
  په بوټانو کښې ښی او کيڼ بېلولے شو،
  نو له هرې جوړې يوه ټاکنه کولے شو:
  ټول ښي بوټان يا ټول کيڼ بوټان اخلو۔
  خو د جرابو دپاره د ټاکنې داسې اصل
  پخپله نه راښکاره کېږي، او تر هغې
  ډاډه کېدلے نه شو چې له هرې جوړې
  يوه جرابه لرونکی ټولګے شته، څو د
  ضرب اصل [يعنې، په عمل کښې انتخاب]
  فرض نه کړو۔
  \citep[p.~126]{Russell1919}
\end{quote}
په لنډه، د لاندې خبرې د ثبوت دپاره
د انتخاب يو ډول ته اړتيا ده: که تاسو
د شمېر وړ ډېرې د جرابو جوړې لرئ،
نو يوازې د شمېر وړ ډېرې جرابې لرئ۔
په حقيقت کښې، بې له د شمېر وړ انتخابه
(يا يې له کوم همارز اصل څخه) د شمېر وړ
سټونو د شمېر وړ اتحاد کېدے شي د شمېر
وړ نه وي۔
\end{ex}

پايله دا ده چې د شمېر وړ انتخاب ډېری
ځله بې له دې کارول کېده چې کاروونکي
يې پوره خبر وي۔ فلسفي پوښتنه دا ده:
د شمېر وړ انتخاب څنګه \emph{توجيه}
کولے شو؟

د وجداني توجيې يوه هڅه ښايي
له فوق‌العمل څخه مرسته وغواړي۔ فرض
کړئ لومړۍ ټاکنه په $\nicefrac{1}{2}$
دقيقه کښې، دويمه په
$\nicefrac{1}{4}$ دقيقه کښې، \dots،
$n$مه ټاکنه په
$\nicefrac{1}{2^n}$ دقيقه کښې،
\dots\@ کوو۔ نو په $1$~دقيقه کښې
به د انتخابونو يوه $\omega$-لړۍ
ترسره کړو او د انتخاب تابعه به
تعريف کړو۔

خو دا ډول فکري تجربه په رښتيا څه
راښودلے شي؟ لومړے دا چې د «ټاکلو»
فکر تر ډېره په لفظي مانا اخلي۔
بل دا چې رياضي له مابعدالطبيعي
امکان سره تړي۔

تر دې مهمه خبره دا ده چې دا د
\emph{مطلق} انتخاب توجيه نه شي
کولے؛ يوازې \emph{د شمېر وړ} انتخاب
ته رسېږي۔ ځکه که \emph{هر} سټ
د انتخاب تابعه غواړي، نو د «هرې
ترتيبي اوږدوالي فوق‌العمل» ترسره
کولو توان به هم پکار وي۔ په ډاګه
ووايو، دا فکر د خندا وړ دے۔

\end{document}